% ============================================================
\chapter{LP Duality}
\label{ch:duality}
% ============================================================

Every optimisation problem hides a mirror: its \emph{dual}. This idea, which traces back to John von Neumann and his work on game theory in the 1940s, is one of the deepest in all of optimisation. In linear programming, duality takes a particularly elegant form: to every minimisation problem corresponds a maximisation problem whose optimal value coincides with that of the primal (the strong duality theorem). The dual provides not only certified bounds, but also an economic interpretation: dual variables are ``shadow prices'' measuring the marginal value of each constraint.

\section{Motivation}

Consider a minimization problem.  We seek to bound the optimal value:
\begin{itemize}
  \item Every feasible solution provides an \textbf{upper bound}.
  \item Can we systematically obtain \textbf{lower bounds}?
\end{itemize}
Duality theory answers affirmatively by associating to each LP (the
\textbf{primal}) another LP (the \textbf{dual}).

\section{Constructing the dual}

\begin{definition}[Primal and dual]
\label{def:primal_dual}
Let the primal (P) be in canonical form:
\[
\text{(P)} \quad
\begin{aligned}
  \max \quad & \mathbf{c}^\top \mathbf{x} \\
  \text{s.t.} \quad & A\mathbf{x} \le \mathbf{b} \\
  & \mathbf{x} \ge \mathbf{0}
\end{aligned}
\]
The \textbf{dual} (D) is:
\[
\text{(D)} \quad
\begin{aligned}
  \min \quad & \mathbf{b}^\top \mathbf{y} \\
  \text{s.t.} \quad & A^\top \mathbf{y} \ge \mathbf{c} \\
  & \mathbf{y} \ge \mathbf{0}
\end{aligned}
\]
where $\mathbf{y} \in \R^m$ is the vector of \textbf{dual variables}.
\end{definition}

\begin{keyformulas}[title=\textbf{Dual construction rules}]
\begin{center}
\renewcommand{\arraystretch}{1.3}
\begin{tabular}{lll}
  \toprule
  \textbf{Primal (max)} & & \textbf{Dual (min)} \\
  \midrule
  Constraint $\le$ & $\leftrightarrow$ & Variable $y_i \ge 0$ \\
  Constraint $\ge$ & $\leftrightarrow$ & Variable $y_i \le 0$ \\
  Constraint $=$ & $\leftrightarrow$ & Variable $y_i$ free \\
  Variable $x_j \ge 0$ & $\leftrightarrow$ & Constraint $\ge$ \\
  Variable $x_j \le 0$ & $\leftrightarrow$ & Constraint $\le$ \\
  Variable $x_j$ free & $\leftrightarrow$ & Constraint $=$ \\
  \bottomrule
\end{tabular}
\end{center}
\end{keyformulas}

\begin{remark}
The dual of the dual is the primal: $(D(D)) = (P)$.
\end{remark}

\section{Duality theorems}

\begin{theorem}[Weak duality]
\label{thm:weak_duality}
If $\mathbf{x}$ is feasible for (P) and $\mathbf{y}$ is feasible for (D),
then:
\[
  \mathbf{c}^\top \mathbf{x} \le \mathbf{b}^\top \mathbf{y}
\]
\end{theorem}

\begin{proof}
We have $A\mathbf{x} \le \mathbf{b}$ and $\mathbf{y} \ge \mathbf{0}$, so
$\mathbf{y}^\top A\mathbf{x} \le \mathbf{y}^\top \mathbf{b}$.  Similarly,
$A^\top\mathbf{y} \ge \mathbf{c}$ and $\mathbf{x} \ge \mathbf{0}$ give
$\mathbf{x}^\top A^\top\mathbf{y} \ge \mathbf{c}^\top\mathbf{x}$.  Hence
$\mathbf{c}^\top\mathbf{x} \le \mathbf{y}^\top A\mathbf{x} \le
\mathbf{b}^\top\mathbf{y}$.
\end{proof}

\begin{corollary}[Bounds]
\begin{enumerate}
  \item If (P) is unbounded ($\max = +\infty$), then (D) is infeasible.
  \item If (D) is unbounded ($\min = -\infty$), then (P) is infeasible.
\end{enumerate}
\end{corollary}

\begin{theorem}[Strong duality]
\label{thm:strong_duality}
If (P) has a finite optimal solution $\mathbf{x}^*$, then (D) also has a
finite optimal solution $\mathbf{y}^*$ and:
\[
  \mathbf{c}^\top \mathbf{x}^* = \mathbf{b}^\top \mathbf{y}^*
\]
\end{theorem}

\begin{intuition}[title=\textbf{Interpretation}]
Strong duality states that the best lower bound (from the dual) coincides
with the best upper bound (from the primal) at optimality.  There is ``no
gap'' between the two.
\end{intuition}

\section{Complementary slackness}

\begin{theorem}[Complementary slackness]
\label{thm:complementary_slackness}
Let $\mathbf{x}^*$ and $\mathbf{y}^*$ be feasible solutions for the primal
and dual respectively.  They are both optimal if and only if:
\[
  y_i^* \left(b_i - \sum_{j} a_{ij} x_j^*\right) = 0, \quad i = 1, \ldots, m
\]
\[
  x_j^* \left(\sum_{i} a_{ij} y_i^* - c_j\right) = 0, \quad j = 1, \ldots, n
\]
\end{theorem}

In other words:
\begin{itemize}
  \item If a primal constraint is \textbf{not tight}, the corresponding dual
        variable is zero.
  \item If a primal variable is \textbf{strictly positive}, the corresponding
        dual constraint is tight.
\end{itemize}

\section{Economic interpretation}

\begin{definition}[Shadow price]
The dual variable $y_i^*$ associated with the $i$-th primal constraint is
called the \textbf{shadow price} of resource $i$.  It represents the
marginal change in the optimal value when the right-hand side $b_i$
increases by one unit:
\[
  y_i^* = \frac{\partial z^*}{\partial b_i}
\]
\end{definition}

\begin{example}[Shadow price interpretation]
Consider the production problem from the previous chapter:
\[
\begin{aligned}
  \max \quad & z = 5x_1 + 4x_2 \\
  \text{s.t.} \quad
    & 6x_1 + 4x_2 \le 24 \quad (y_1) \\
    & x_1 + 2x_2 \le 6 \quad (y_2) \\
    & x_1, x_2 \ge 0
\end{aligned}
\]

The dual is:
\[
\begin{aligned}
  \min \quad & w = 24y_1 + 6y_2 \\
  \text{s.t.} \quad
    & 6y_1 + y_2 \ge 5 \\
    & 4y_1 + 2y_2 \ge 4 \\
    & y_1, y_2 \ge 0
\end{aligned}
\]

From the final simplex tableau (Chapter~3):
$y_1^* = 3/4$, $y_2^* = 1/2$.  Check: $w^* = 24(3/4) + 6(1/2) = 21 = z^*$.

Interpretation: increasing the capacity of constraint~1 from 24 to 25
increases $z^*$ by approximately $3/4 = 0.75$.
\end{example}

\section{Reading the dual from the simplex tableau}

\begin{proposition}[Dual variables from the final tableau]
In the final simplex tableau (primal in standard form with slack variables),
the optimal dual variable values can be read from the \textbf{objective row}
at the columns of the slack variables:
\[
  y_i^* = \bar{c}_{s_i} \quad \text{(reduced cost of slack variable $s_i$)}
\]
\end{proposition}

\section{The dual simplex method}

The \textbf{dual simplex} method solves the dual directly on the primal
tableau.  It is useful when the basis is \textbf{dual feasible} (all
reduced costs $\ge 0$) but not primal feasible ($\bar{b}_i < 0$ for some
$i$).

\begin{algorithme}[title=\textbf{Dual simplex --- one iteration}]
\begin{enumerate}
  \item \textbf{Primal optimality test}: if $\bar{b}_i \ge 0$ for all $i$,
        \textsc{stop} (optimal solution).
  \item \textbf{Leaving variable}: choose $r = \argmin_i \bar{b}_i$ (most
        negative row).
  \item \textbf{Infeasibility test}: if $\bar{a}_{rj} \ge 0$ for all $j$,
        \textsc{stop} (primal is infeasible).
  \item \textbf{Entering variable}: apply the dual ratio:
        \[
          k = \argmin_{j : \bar{a}_{rj} < 0}
            \frac{\bar{c}_j}{\abs{\bar{a}_{rj}}}
        \]
  \item \textbf{Pivot}: pivot on $\bar{a}_{rk}$.
\end{enumerate}
\end{algorithme}

\section{Numerical example}

\begin{example}[Constructing and solving the dual]
\[
\text{(P)} \quad
\begin{aligned}
  \max \quad & z = 4x_1 + 3x_2 \\
  \text{s.t.} \quad
    & 2x_1 + x_2 \le 10 \\
    & x_1 + 3x_2 \le 12 \\
    & x_1, x_2 \ge 0
\end{aligned}
\]

The dual:
\[
\text{(D)} \quad
\begin{aligned}
  \min \quad & w = 10y_1 + 12y_2 \\
  \text{s.t.} \quad
    & 2y_1 + y_2 \ge 4 \\
    & y_1 + 3y_2 \ge 3 \\
    & y_1, y_2 \ge 0
\end{aligned}
\]
\end{example}

\begin{codeblock}[title=\textbf{Verification with PuLP}]
\begin{minted}{python}
import pulp

# Primal
P = pulp.LpProblem("Primal", pulp.LpMaximize)
x1 = pulp.LpVariable("x1", lowBound=0)
x2 = pulp.LpVariable("x2", lowBound=0)
P += 4*x1 + 3*x2
P += 2*x1 + x2 <= 10, "c1"
P += x1 + 3*x2 <= 12, "c2"
P.solve(pulp.PULP_CBC_CMD(msg=0))
print(f"Primal: z* = {pulp.value(P.objective):.2f}")
print(f"  x* = ({x1.varValue}, {x2.varValue})")

# Dual
D = pulp.LpProblem("Dual", pulp.LpMinimize)
y1 = pulp.LpVariable("y1", lowBound=0)
y2 = pulp.LpVariable("y2", lowBound=0)
D += 10*y1 + 12*y2
D += 2*y1 + y2 >= 4
D += y1 + 3*y2 >= 3
D.solve(pulp.PULP_CBC_CMD(msg=0))
print(f"Dual: w* = {pulp.value(D.objective):.2f}")
print(f"  y* = ({y1.varValue}, {y2.varValue})")
\end{minted}
\end{codeblock}

\begin{codeoutput}
Primal: z* = 22.80
  x* = (5.4, 2.2)
Dual: w* = 22.80
  y* = (1.8, 0.4)
\end{codeoutput}

\section{Summary of primal-dual relationships}

\begin{center}
\renewcommand{\arraystretch}{1.3}
\begin{tabular}{lccc}
  \toprule
  & \textbf{(P) feasible} & \textbf{(P) infeasible} \\
  \midrule
  \textbf{(D) feasible} & Equal finite optima & (D) unbounded \\
  \textbf{(D) infeasible} & (P) unbounded & Both infeasible \\
  \bottomrule
\end{tabular}
\end{center}

\begin{remark}
It is possible for both primal and dual to be infeasible.  Example:
$\max\ x_1 - x_2$ s.t.\ $-x_1 + x_2 \le -1$, $x_1 - x_2 \le -1$,
$x_1, x_2 \ge 0$.
\end{remark}

\section{Exercises}

\begin{exercise}[$\star$ --- Write the dual]
Write the dual of the following LPs:
\begin{enumerate}
  \item $\max\ 3x_1 + 2x_2$ s.t.\ $x_1 + x_2 \le 5$,
        $2x_1 + x_2 \le 8$, $x_1, x_2 \ge 0$.
  \item $\min\ 5x_1 + 3x_2$ s.t.\ $x_1 + x_2 = 4$,
        $2x_1 + x_2 \ge 6$, $x_1, x_2 \ge 0$.
\end{enumerate}
\end{exercise}

\begin{exercise}[$\star\star$ --- Strong duality verification]
Solve both the primal and dual from Exercise~1(a) and verify that the optimal
values are equal.
\end{exercise}

\begin{exercise}[$\star\star$ --- Complementary slackness]
Use complementary slackness to find the optimal dual solution given the
primal solution $x_1^* = 3$, $x_2^* = 2$ for the LP:
$\max\ 4x_1 + 3x_2$ s.t.\ $2x_1 + x_2 \le 8$, $x_1 + 2x_2 \le 7$,
$x_1, x_2 \ge 0$.
\end{exercise}

\begin{exercise}[$\star\star$ --- Economic interpretation]
A factory has three resources $R_1, R_2, R_3$ with capacities 100, 80, and
120.  The optimal shadow prices are $y_1^* = 2$, $y_2^* = 0$, $y_3^* = 3$.
Interpret these values.  A supplier offers 10 additional units of $R_2$ for
\$5.  Should the offer be accepted?
\end{exercise}

\begin{exercise}[$\star\star\star$ --- Proof]
Prove the weak duality theorem for the general case (standard form primal).
\end{exercise}

\begin{exercise}[$\star\star\star$ --- Dual simplex]
Apply the dual simplex to the problem:
\[
\begin{aligned}
  \min \quad & z = 2x_1 + 3x_2 + x_3 \\
  \text{s.t.} \quad
    & x_1 + x_2 + x_3 \ge 6 \\
    & 2x_1 + x_3 \ge 4 \\
    & x_1, x_2, x_3 \ge 0
\end{aligned}
\]
\end{exercise}
